\subsection{Extension and restriction of the structure group}

6.6.1. Let \(\lambda = (\mathbf{P}, \mathbf{G}, \mathbf{B}, \pi)\) be a principal fibration, and let \(\varphi\) be a homomorphism from \(\mathbf{G}\) into a group variety \(\mathbf{H}\) . Let us make \(\mathbf{G}\) operate on the left on \(\mathbf{H}\) by \(g \cdot h = \varphi(g) \cdot h\) and let \(\mathbf{P} \times^{\mathbf{G}} \mathbf{H}\) be the associated fiber bundle to \(\lambda\) with typical fiber \(\mathbf{H}\) . Since the right translations in \(\mathbf{H}\) are compatible with the operations of \(\mathbf{G}\) , the group \(\mathbf{H}\) operates on the right on \(\mathbf{P} \times^{\mathbf{G}} \mathbf{H}\) (cf.~no. 6.5.5); if \(\pi_{\mathbf{H}}\) denotes the projection of \(\mathbf{P} \times^{\mathbf{G}} \mathbf{H}\) onto \(\mathbf{B}\) , the quadruple ( \(\mathbf{P} \times^{\mathbf{G}} \mathbf{H}\) , \(\mathbf{H}\) , \(\mathbf{B}\) , \(\pi_{\mathbf{H}}\) ) is a principal fibration, denoted \(\varphi(\lambda)\) ; it is said to be deduced from \(\lambda\) by the homomorphism \(\varphi\) .

The map \(f\) of \(\mathbf{P}\) to \(\mathbf{P} \times^{\mathbf{G}}\mathbf{H}\) which associates to \(x \in \mathbf{P}\) the class of \((x, e)\) is a \(\mathbf{B}\) -morphism of \(\mathbf{P}\) to \(\mathbf{P} \times^{\mathbf{G}}\mathbf{H}\) compatible with \(\varphi\) (cf.~no. 6.3.1). Moreover, if \(f'\) is a \(\mathbf{B}\) -morphism of \(\mathbf{P}\) into a principal fiber bundle \(\mathbf{P}'\) of the structural group \(\mathbf{H}\) and if \(f'\) is compatible with \(\varphi\) , there exists a unique \(\mathbf{H}\) - \(\mathbf{B}\) -isomorphism \(\theta\) of \(\mathbf{P} \times^{\mathbf{G}}\mathbf{H}\) onto \(\mathbf{P}'\) such that \(f' = \theta \circ f\) .

6.6.2. Suppose that \(\lambda\) is defined by means of an open covering \(\mathcal{V} = (\mathrm{U}_i)\) of B and of a cocycle \((g_{ij})\) (6.4.2). Then \(\varphi(\lambda)\) can be defined using the same covering and the cocycle \((h_{ij})\) , with \(h_{ij} = \varphi \circ g_{ij}\) .

6.6.3. Let F be a variety on which the group H acts on the left; we denote by \((h, y) \mapsto h \cdot y\) the action law of H on F. The group G acts on F by \((g, y) \mapsto \varphi(g) \cdot y\) . Let E be a fiber bundle associated with \(\varphi(\lambda)\) with typical fiber F. The map \((x, y) \mapsto f(x) \cdot y\) from P × F to E (the map f being that defined in no. 6.6.1) endows E with a structure of fiber bundle associated with \(\lambda\) of fiber type F. In particular, \((\mathbf{P} \times^{\mathbf{G}}\mathbf{H}) \times^{\mathbf{H}}\mathbf{F}\) is identified with P × \(^{\mathbf{G}}\) F.

6.6.4. Suppose that G is a Lie subgroup of H, and that \(\varphi: \mathbf{G} \to \mathbf{H}\) is the canonical injection of G into H. One then says that \(\varphi(\lambda)\) is deduced from \(\lambda\) by extension to H of the structural group. The morphism \(f: \mathbf{P} \to \mathbf{P} \times^{\mathbf{G}} \mathbf{H}\) of No. 6.6.1 is an isomorphism from P onto a closed submanifold of \(\mathbf{P} \times^{\mathbf{G}} \mathbf{H}\) (equal to \(\mathbf{P} \times^{\mathbf{G}} \mathbf{H}\) if \(\mathbf{G} = \mathbf{H}\) ); this isomorphism is compatible with the operations of G (note that G acts on \(\mathbf{P} \times^{\mathbf{G}} \mathbf{H}\) as a subgroup of H).

6.6.5. Suppose further that G is a Lie subgroup of H, and let \(\mu = (Q, H, B, \pi)\) be a principal fibration with structure group H and base B, and let E be a fiber bundle associated with \(\mu\) of fiber type H/G. If \(\gamma\) denotes the canonical projection of H onto H/G, set \(\delta(y) = y \cdot \gamma(e)\) , for \(y \in Q\) (where \(e\) denotes the identity element of H); one thus obtains a morphism \(\delta: Q \to E\) and the quadruple \((Q, G, E, \delta)\) is a principal fibration. In particular, Q/G is canonically identified with E, itself isomorphic to \(Q \times^{H} H/G^{1}\) .

Now let \(s: \mathbf{B} \to \mathbf{E}\) be a section of \(\mathbf{E}\) , and let \(\lambda_s\) be the inverse image under \(s\) of the fibration \((\mathbf{Q}, \mathbf{G}, \mathbf{E}, \delta)\) that we have just defined. It is a principal fibration, with base \(\mathbf{B}\) and with structural group \(\mathbf{G}\) , and its extension to \(\mathbf{H}\) is isomorphic to \(\mu\) ; any principal fibration enjoying these properties can be obtained in this way, up to isomorphism; two sections \(s_1\) and \(s_2\) of \(\mathbf{E}\) define isomorphic fibrations if and only if they are transformed into one another by an H-B-automorphism of \(\mu\) .

